\begin{abcliste}
\abc The magnetic force is the centripetal force, $e v B = m_p v^2/r$, so $r = \frac{m_p v}{e B}$. The period of one turn is
\begin{align}
 T &= \frac{2\pi r}{v} = \frac{2\pi m_p}{e B}
\end{align}
It does not depend on the speed: a voltage of constant frequency stays in step with the protons. The frequency is
\begin{align}
 f &= \frac{1}{T} = \fcF \\
   &= \frac{\nce\times\B}{2\pi\times\ncmp} \\
   &= \fc \approx \result{\fcP{6}{2}}
\end{align}
(The angular frequency $\omega = 2\pi f$ or twice the frequency would be the typical mistakes.)

\abc At the edge, $r = R$:
\begin{align}
 v &= \vF \\
   &= \frac{\nce\times\B\times\R}{\ncmp} \\
   &= \v \approx \result{\vP{0}{2}}
\end{align}
This is about a quarter of the speed of light, so the classical formulas are still a fair approximation. The kinetic energy is
\begin{align}
 \ssc{E}{kin} &= \frac{1}{2}m_p v^2 = \EkF \\
   &= \frac{\left(\nce\times\B\times\R\right)^2}{2\times\ncmp} \\
   &= \Ek \approx \result{\EkeVP{6}{2}}
\end{align}

\abc The speed $v = \frac{eBR}{m_p}$ grows with $R$, the energy $\frac{1}{2}m_p v^2$ with $R^2$: with twice the radius the final energy is $\result{\text{four times as large}}$.
\end{abcliste}
